\hmsubsection{Object Detection Methods}{OAH}{}

Object detection is a well-established problem in computer vision, and a wide range of
methods have been developed to address it. These methods can broadly be divided into two
categories: classical image processing techniques and learning-based approaches.

Classical methods rely on hand-crafted features and deterministic algorithms. Colour
thresholding in the HSV colour space is a widely used technique for detecting objects of known
colour, as HSV separates chromatic information from illumination, making thresholds more
robust to lighting variations~\cite{bradski2008learning}. The Hough Circle Transform is
another classical method that detects circular shapes by voting in a parameter space defined
by circle centre and radius~\cite{hough1962method}. These methods are computationally
efficient and well-suited for embedded systems with limited processing power, such as the
Raspberry Pi~5 used in this project.

Learning-based approaches, such as convolutional neural networks (CNNs) and object
detectors like YOLO~\cite{redmon2016yolo} and SSD, have demonstrated superior accuracy on
complex detection tasks in uncontrolled environments with many object classes. However, they
typically require significant computational resources and large labelled datasets for training.
During early development of this project, MobileNetV2 transfer learning was evaluated as a
detection method. It produced inference times of 200--500~ms per frame on the Raspberry
Pi~5, overfit on the limited training dataset, and required extensive labelled data for each
new lighting condition. This made it unsuitable for real-time operation.

In this project, a hybrid classical approach is employed instead. HSV colour segmentation
and the Hough Circle Transform serve as the primary detection methods, while a lightweight
Support Vector Machine (SVM) classifier was evaluated as a secondary colour verification
step. The classifier was subsequently disabled after being found to misclassify dark
navy-blue balls as red with high confidence, and the HSV and Hough ensemble alone was
found to be sufficiently reliable. This combination achieves 20--25~FPS on the Raspberry
Pi~5, a detection rate of 98--100\% for both red and blue balls under laboratory conditions,
and a false positive rate below 2\%. The fully deterministic pipeline also offers complete
explainability --- each rejection has a specific geometric or chromatic reason visible in the
diagnostic tools --- which is a significant advantage over black-box deep learning approaches
for a controlled sorting environment.